Due to its long-standing focus on flood prevention, the Netherlands represents a unique case in global water management. With approximately 26 percent of its land below mean sea level, the Dutch water system has historically been optimised for rapid drainage, water discharge, and protection against flooding. This system is increasingly challenged by a changing hydro-climatic reality.

Globally, the frequency and severity of droughts have increased over recent decades, driven primarily by climate change and associated disruptions to the hydrological cycle \cite{Shyrokaya2024Advances}. The Netherlands is no exception to this trend. Severe drought events have occurred in 1976, 2018, 2019, 2020 and 2022 highlighting a growing exposure to prolonged precipitation deficits in a system originally designed for excess water removal rather than water retention \cite{KNMI2023}. 

This shift introduces a structural hydrological trade-off: infrastructure optimised for flood protection may inadvertently lead to more water scarcity during dry periods. Reduced precipitation, higher evapotranspiration, and altered river discharges collectively contribute to increased drought stress across multiple sectors. In Western Europe, climatological assessments indicate both an increase in the frequency and intensity of extreme drought conditions, reinforcing the need for adaptive water management strategies \cite{KNMI2023}.

At the policy level, climate change adaptation is becoming increasingly central. Discussions at the European Climate Change Adaptation Conference (ECCA) 2025 emphasised the growing importance of adaptation strategies alongside mitigation efforts, particularly given the challenges of fully achieving long-term emission targets. Effective drought adaptation frameworks rely on three core components: anticipating drought events, understanding their potential impacts, and designing appropriate response strategies \cite{Bulut2026Framework}.

This thesis contributes to these adaptation frameworks in two ways. First, it develops a drought impact database to improve understanding of the potential consequences of drought events. Second, it develops forecasting models to better anticipate the consequences of future droughts, thereby supporting more proactive drought management.
